% Generated draft: requires mathematical and source review.
\documentclass[11pt]{article}
\usepackage[margin=25mm]{geometry}
\usepackage{fontspec}
\setmainfont{Latin Modern Roman}
\usepackage{amsmath,amssymb,mathtools,unicode-math}
\setmathfont{Latin Modern Math}
\usepackage{enumitem,needspace,xurl,hyperref}
\setlength{\parindent}{0pt}
\setlength{\parskip}{5pt}
\setlength{\emergencystretch}{3em}
\newcommand{\N}{\mathbb{N}}
\newcommand{\Z}{\mathbb{Z}}
\newcommand{\Q}{\mathbb{Q}}
\newcommand{\R}{\mathbb{R}}
\newcommand{\C}{\mathbb{C}}
\newcommand{\F}{\mathbb{F}}
\newcommand{\A}{\mathbb{A}}
\newcommand{\PP}{\mathbb{P}}
\newcommand{\E}{\mathbb{E}}
\newcommand{\eps}{\varepsilon}
\newcommand{\dd}{\,\mathrm{d}}
\begin{document}
\section*{UnsolvedMath ID: 800134}
\textbf{Source problem code:} catalogue entry
\subsection{Definitions and mathematical statement}
For a real number \(x\), irrationality means \(x\in\R\setminus\Q\), where \(\Q\) denotes the rational numbers. The supplied source requests:
\begin{quote}
Show that \(\gamma_1\) is irrational.
\end{quote}
However, the record does not define \(\gamma_1\), specify its domain, or identify the function or other object from which it is obtained. If the intended \(\gamma_1\) is a specified real number, the requested conclusion would be
\[
\gamma_1\notin\Q.
\]
This conditional interpretation does not identify the intended number; the supplied record is insufficient for a self-contained mathematical formulation.

The surrounding source discusses an undefined condition called ``LI'' for higher-degree \(L\)-functions. It cautions about factorizations of Dedekind \(\zeta\)-functions and Artin \(L\)-functions into other \(L\)-functions, vanishing of elliptic curve \(L\)-functions at the central point, and the possibility of shifting an entire \(L\)-function vertically when no restriction to the central character is imposed. It calls the irrationality question a concrete instance of an earlier problem, whose formulation is absent. These contextual qualifications do not supply a definition of \(\gamma_1\) or precise hypotheses for the requested assertion.
\subsection{Short English statement}
The source asks for a proof that \(\gamma_1\) is irrational, but does not identify \(\gamma_1\). Its surrounding discussion concerns qualifications to an undefined ``LI'' condition for higher-degree \(L\)-functions. The intended number and the missing preceding formulation must be recovered before a precise statement can be completed.
\subsection{Sources}
\begin{itemize}
\item Supplied local UnsolvedMath registry record: permanent ID \(800134\), exact source problem code \texttt{catalogue entry}, title ``Comparative Prime Number Theory --- Problem 3.'' The registry labels the record open; this is supplied status metadata, not an independently verified mathematical assessment.
\item Hamieh, Kadiri, Martin, and Ng, \emph{Comparative Prime Number Theory Problem List} (2024), source item Problem (\#3), as identified by the supplied record: \url{https://arxiv.org/abs/2407.03530}. The background records access on 2026-10-16 and extraction by \texttt{arxiv\_latex}; these are supplied provenance details, not a claim of independent consultation or verification.
\item The registry attributes the proposal to Hamieh, Kadiri, Martin, and Ng (2024). The extracted statement also contains the attribution ``Alexandre Bailleul, Universit\'e Paris--Saclay'' before the concrete-instance sentence and ``Nicol Leong, University of New South Wales--Canberra at ADFA'' after the irrationality question. Their precise attachment to the surrounding problems is unclear from this extract.
\end{itemize}
\end{document}
